\appendix
\section*{Appendix}
\subsection*{{\bf A.} Some Examples}
\label{sec:examples}
The first example shows that the factor $1/\sqrt{m}$ in the statement of $\rsh$ (\Cref{rsh-0}) cannot be improved, even for a simple polytope $K$ consisting of just one line segment. 

\vspace*{0.1 in}
\noindent
{\bf Example 1:} Suppose  $K$  is a polytope in $\Re^d$ with just two vertices $\Delta$ distance apart and $V = \Re^d, \delta=1$. It is easy to verify from standard results on random projection on a line that the additive term in inequality~\eqref{eq:rsh} is tight up to constant factors.

\vspace*{0.1 in}
The next example shows that the success probability of the desired event~\eqref{eq:rsh} in~\Cref{rsh-0} needs to depend on $k$, the number of vertices of the polytope $K$. 

\vspace*{0.1 in}
\noindent
{\bf Example 2:} Consider the Euclidean space $\Re^d$ and the $d-1$ dimensional sphere $K:=\{x \in \Re^d: |x|\leq 1, x_1=0\}$. Observe that the diameter $\dia$ of $K$ is 1. 
Let $a$ denote the point $(1,0,0,\ldots ,0)$. It is at distance 1 from $K$, and so the parameter $\delta$ in the statement of~\Cref{rsh-0} can be set to 1.  But, with high probability,  a random unit length vector $u$ has $u\cdot a\approx \frac{c}{\sqrt d}$ for some constant $c$,  whereas the maximum of  $v\cdot x$ over all points $x$ in the sphere $K$ is about 1. Therefore the probability of the event~\eqref{eq:rsh} is exponentially small in $d$.